\section{Kimi K2 y la capacidad de agente}

A partir del optimizador Muon y de técnicas de entrenamiento estable, Kimi K2 logró varios avances tras el post-entrenamiento con aprendizaje por refuerzo:

\begin{itemize}
\item \textbf{Mejora integral de la capacidad de agente}: puede compararse con las empresas de vanguardia de Estados Unidos
\item \textbf{HLE (Humanities Last Exam)}: alcanza 45\% de exactitud, por encima de OpenAI; este es un benchmark de alta dificultad que incluso a los humanos les cuesta trabajo resolver
\item \textbf{K2 Thinking}: completa entre 200 y 300 pasos consecutivos de llamadas a herramientas, con razonamiento, búsqueda y escritura de programas en Python de manera continua durante el proceso
\item \textbf{Kimi K2 es el primer modelo agéntico de China}
\end{itemize}

\begin{knowledgebox}{Por qué un agente es, en esencia, un problema de búsqueda}
La inferencia o el entrenamiento con RL de un agente es, en esencia, un proceso de búsqueda. Por ejemplo, desarrollar desde cero un sistema operativo Linux: si se dispusiera de cómputo infinito, se podrían enumerar todas las combinaciones de tokens hasta encontrar la respuesta correcta. Un modelo base más fuerte ofrece un prior (heurística) más sólido, lo que reduce el espacio de búsqueda. Un aumento en la eficiencia de tokens equivale a un prior más fuerte, mientras que el contexto largo ofrece una memoria de trabajo más sólida (capacidad de percepción del entorno).
\end{knowledgebox}

